\documentclass[QCM, noCustomPackages]{UPSTI_Document}
\usepackage{IUT_Annecy}
\usepackage{logos_iut}
% \QCMchemin / \QCMcheminModule permettent de recompiler ce fichier depuis
% un autre répertoire (ex. déploiement AMC via un fiche_evaluation.tex qui
% redéfinit ces macros en chemins absolus avant de \input ce fichier) : un
% \newcommand{\sequence}{2}
\newcommand{\UPSTIsousTitreEnTete}{Sequence \sequence~: Boucles et fonctions}


\graphicspath{{images/}}
 littéral se résout par rapport au répertoire de
% travail du processus LaTeX, pas par rapport à ce fichier — donc pas
% "recompilable ailleurs" tel quel. Valeur par défaut inchangée : compiler
% ce fichier tel quel (ex. depuis supports-informatique/) fonctionne
% exactement comme avant.
\providecommand{\QCMchemin}{..}
\providecommand{\QCMcheminModule}{../..}
\input{\QCMchemin/preamble.tex}
\input{\QCMcheminModule/preamble_module.tex}

\def\photodir{photos}


\documentVersion{E}
\newcommand{\UPSTInumeroVersion}{0}
\newcommand{\UPSTImessage}{test}

\newcommand{\UPSTItitre}{QCM --- Variables, conditions et boucles}


\newcommand{\sujet}{% Début de la boucle sur les étudiants
\exemplaire{1}{%
\renewcommand{\UPSTImessage}{\champnom{\nom~\prenom}}
\renewcommand{\UPSTInumeroVersion}{\id{}}
% \photo n'est défini par \csvreader[head to column names,] que si le CSV
% étudiants a réellement une colonne "photo" -- sinon \photo n'existe pas
% du tout (pas juste vide) et cette ligne casserait le logo d'en-tête pour
% CHAQUE étudiant ("Undefined control sequence" + image introuvable),
% comme observé avec etudiants-GEII-S1.csv (pas de colonne photo). On ne
% remplace donc le logo IUT par défaut (logos_iut.sty) que si la colonne
% existe pour ce CSV.
\ifdefined\photo\renewcommand{\UPSTIlogoLycee}{\photo}\fi
\UPSTIbuildPage
\setcounter{figure}{0}
\setcounter{table}{0}

\section{Variables simples}
\restituegroupe[3]{variables}
\section{Affectations et expressions}
\restituegroupe[3]{affectations}
\section{Conditions}
\restituegroupe[6]{conditions}
\section{Boucles simples}
\restituegroupe[5]{boucles}
\section{Boucles imbriquées}
\restituegroupe[3]{boucles-imbriquees}
\AMClabel{UPSTIfinexemplaire}
\AMCcleardoublepage
}
}




\begin{document}
\setdefaultgroupmode{withoutreplacement}
\AMCrandomseed{1237893}
% ==================== PREAMBULE LOCAL — QCM variables-boucles ====================
% Variante locale de \evaluationProfGrille remplaçant le quadrillage \casesGrid
% par un tableau structuré (trace d'exécution), sans toucher au paquet UPSTI
% partagé. Le mécanisme de notation (\bareme / \AMCOpen via \evaluationProfBase)
% est strictement identique : seul l'habillage visuel (lineuptext) change.
%
% Usage : \evaluationProfTableauTrace{<bareme max>}{<contenu LaTeX du tableau>}
% (argument optionnel #1 identique à \evaluationProfGrille, non utilisé ici)
\newcommandx{\evaluationProfTableauTrace}[3][1=]{{%
  \evaluationProfBase[#2]{\max*GRAD/5}{\evaluationGraduee}{%
    lineup=true, lineuptext={#3}, lines=0%
  }%
}}

% ==================== THEME 1 — VARIABLES SIMPLES ====================

% ---- var-declaration ----
\element{bank-var-declaration}{
    \begin{question}{var-declaration-v1}
        Quelle ligne déclare correctement une variable entière nommée \texttt{age} ?
        \begin{multicols}{4}
            \begin{reponses}
                \bonne{\lstinline[language=C]|int age;|}
                \mauvaise{\lstinline[language=C]|int age|}
                \mauvaise{\lstinline[language=C]|age int;|}
                \mauvaise{\lstinline[language=C]|integer age;|}
            \end{reponses}
        \end{multicols}
    \end{question}
}
\element{bank-var-declaration}{
    \begin{question}{var-declaration-v2}
        Quelle ligne déclare correctement une variable entière nommée \texttt{score} ?
        \begin{multicols}{4}
            \begin{reponses}
                \bonne{\lstinline[language=C]|int score;|}
                \mauvaise{\lstinline[language=C]|int score|}
                \mauvaise{\lstinline[language=C]|score int;|}
                \mauvaise{\lstinline[language=C]|integer score;|}
            \end{reponses}
        \end{multicols}
    \end{question}
}
\element{bank-var-declaration}{
    \begin{question}{var-declaration-v3}
        Quelle ligne déclare correctement une variable entière nommée \texttt{compteur} ?
        \begin{multicols}{4}
            \begin{reponses}
                \bonne{\lstinline[language=C]|int compteur;|}
                \mauvaise{\lstinline[language=C]|int compteur|}
                \mauvaise{\lstinline[language=C]|compteur int;|}
                \mauvaise{\lstinline[language=C]|integer compteur;|}
            \end{reponses}
        \end{multicols}
    \end{question}
}

\element{variables}{
    \restituegroupe[1]{bank-var-declaration}
}

% ---- var-type-decimal ----
\element{bank-var-type-decimal}{
    \begin{question}{var-type-decimal-v1}
        Pour stocker un nombre à virgule comme \texttt{3.14}, quel type utiliser en priorité ?
        \begin{multicols}{4}
            \begin{reponses}
                \mauvaise{\lstinline[language=C]|int|}
                \bonne{\lstinline[language=C]|float|}
                \mauvaise{\lstinline[language=C]|char|}
                \mauvaise{\lstinline[language=C]|string|}
            \end{reponses}
        \end{multicols}
    \end{question}
}
\element{bank-var-type-decimal}{
    \begin{question}{var-type-decimal-v2}
        Pour stocker un nombre à virgule comme \texttt{2.5}, quel type utiliser en priorité ?
        \begin{multicols}{4}
            \begin{reponses}
                \mauvaise{\lstinline[language=C]|int|}
                \bonne{\lstinline[language=C]|float|}
                \mauvaise{\lstinline[language=C]|char|}
                \mauvaise{\lstinline[language=C]|string|}
            \end{reponses}
        \end{multicols}
    \end{question}
}
\element{bank-var-type-decimal}{
    \begin{question}{var-type-decimal-v3}
        Pour stocker un nombre à virgule comme \texttt{9.99}, quel type utiliser en priorité ?
        \begin{multicols}{4}
            \begin{reponses}
                \mauvaise{\lstinline[language=C]|int|}
                \bonne{\lstinline[language=C]|float|}
                \mauvaise{\lstinline[language=C]|char|}
                \mauvaise{\lstinline[language=C]|string|}
            \end{reponses}
        \end{multicols}
    \end{question}
}

\element{variables}{
    \restituegroupe[1]{bank-var-type-decimal}
}

% ---- var-type-caractere ----
\element{bank-var-type-caractere}{
    \begin{question}{var-type-caractere-v1}
        Quel type permet de stocker un unique caractère, comme \texttt{'A'} ?
        \begin{multicols}{4}
            \begin{reponses}
                \mauvaise{\lstinline[language=C]|int|}
                \mauvaise{\lstinline[language=C]|float|}
                \bonne{\lstinline[language=C]|char|}
                \mauvaise{\lstinline[language=C]|void|}
            \end{reponses}
        \end{multicols}
    \end{question}
}
\element{bank-var-type-caractere}{
    \begin{question}{var-type-caractere-v2}
        Quel type permet de stocker un unique caractère, comme \texttt{'Z'} ?
        \begin{multicols}{4}
            \begin{reponses}
                \mauvaise{\lstinline[language=C]|int|}
                \mauvaise{\lstinline[language=C]|float|}
                \bonne{\lstinline[language=C]|char|}
                \mauvaise{\lstinline[language=C]|void|}
            \end{reponses}
        \end{multicols}
    \end{question}
}
\element{bank-var-type-caractere}{
    \begin{question}{var-type-caractere-v3}
        Quel type permet de stocker un unique caractère, comme \texttt{'x'} ?
        \begin{multicols}{4}
            \begin{reponses}
                \mauvaise{\lstinline[language=C]|int|}
                \mauvaise{\lstinline[language=C]|float|}
                \bonne{\lstinline[language=C]|char|}
                \mauvaise{\lstinline[language=C]|void|}
            \end{reponses}
        \end{multicols}
    \end{question}
}

\element{variables}{
    \restituegroupe[1]{bank-var-type-caractere}
}

% ---- var-nommage ----
% Durci : le piège "commence par un chiffre" est repéré instantanément par tous
% les débutants ; on le remplace par le piège plus subtil du mot-clé réservé
% utilisé comme nom de variable, et on ajoute un identifiant à underscore
% initial (valide, mais souvent cru interdit) pour forcer une vraie analyse.
\element{bank-var-nommage}{
    \begin{question}{var-nommage-v1}
        Quel nom de variable est \textbf{invalide} en C ?
        \begin{multicols}{5}
            \begin{reponses}
                \mauvaise{\lstinline[language=C]|nb_essais|}
                \mauvaise{\lstinline[language=C]|_compteur|}
                \bonne{\lstinline[language=C]|for|}
                \mauvaise{\lstinline[language=C]|essai2|}
                \mauvaise{\lstinline[language=C]|Essai_Final|}
            \end{reponses}
        \end{multicols}
    \end{question}
}
\element{bank-var-nommage}{
    \begin{question}{var-nommage-v2}
        Quel nom de variable est \textbf{invalide} en C ?
        \begin{multicols}{5}
            \begin{reponses}
                \mauvaise{\lstinline[language=C]|nb_tours|}
                \mauvaise{\lstinline[language=C]|_limite|}
                \bonne{\lstinline[language=C]|while|}
                \mauvaise{\lstinline[language=C]|tour3|}
                \mauvaise{\lstinline[language=C]|TourMax|}
            \end{reponses}
        \end{multicols}
    \end{question}
}
\element{bank-var-nommage}{
    \begin{question}{var-nommage-v3}
        Quel nom de variable est \textbf{invalide} en C ?
        \begin{multicols}{5}
            \begin{reponses}
                \mauvaise{\lstinline[language=C]|note_max|}
                \mauvaise{\lstinline[language=C]|_seuil|}
                \bonne{\lstinline[language=C]|if|}
                \mauvaise{\lstinline[language=C]|note1|}
                \mauvaise{\lstinline[language=C]|NoteFinale|}
            \end{reponses}
        \end{multicols}
    \end{question}
}

\element{variables}{
    \restituegroupe[1]{bank-var-nommage}
}

% ---- var-initialisation ----
\element{bank-var-initialisation}{
    \begin{question}{var-initialisation-v1}
        Que vaut \texttt{x} juste après l'exécution de la ligne \lstinline[language=C]|int x;| (sans affectation) ?
        \begin{multicols}{4}
            \begin{reponses}
                \mauvaise{0}
                \bonne{Une valeur indéterminée}
                \mauvaise{-1}
                \mauvaise{Une erreur de compilation}
            \end{reponses}
        \end{multicols}
    \end{question}
}
\element{bank-var-initialisation}{
    \begin{question}{var-initialisation-v2}
        Que vaut \texttt{n} juste après l'exécution de la ligne \lstinline[language=C]|int n;| (sans affectation) ?
        \begin{multicols}{4}
            \begin{reponses}
                \mauvaise{0}
                \bonne{Une valeur indéterminée}
                \mauvaise{-1}
                \mauvaise{Une erreur de compilation}
            \end{reponses}
        \end{multicols}
    \end{question}
}
\element{bank-var-initialisation}{
    \begin{question}{var-initialisation-v3}
        Que vaut \texttt{total} juste après l'exécution de la ligne \lstinline[language=C]|int total;| (sans affectation) ?
        \begin{multicols}{4}
            \begin{reponses}
                \mauvaise{0}
                \bonne{Une valeur indéterminée}
                \mauvaise{-1}
                \mauvaise{Une erreur de compilation}
            \end{reponses}
        \end{multicols}
    \end{question}
}

\element{variables}{
    \restituegroupe[1]{bank-var-initialisation}
}

% ==================== THEME 2 — AFFECTATIONS / EXPRESSIONS ====================

% ---- aff-egal-vs-double-egal ----
\element{bank-aff-egal-vs-double-egal}{
    \begin{question}{aff-egal-vs-double-egal-v1}
        Pour affecter la valeur \texttt{10} à la variable \texttt{x}, quelle écriture est correcte ?
        \begin{multicols}{4}
            \begin{reponses}
                \bonne{\lstinline[language=C]|x = 10;|}
                \mauvaise{\lstinline[language=C]|x == 10;|}
                \mauvaise{\lstinline[language=C]|10 = x;|}
                \mauvaise{\lstinline[language=C]|x := 10;|}
            \end{reponses}
        \end{multicols}
    \end{question}
}
\element{bank-aff-egal-vs-double-egal}{
    \begin{question}{aff-egal-vs-double-egal-v2}
        Pour affecter la valeur \texttt{25} à la variable \texttt{n}, quelle écriture est correcte ?
        \begin{multicols}{4}
            \begin{reponses}
                \bonne{\lstinline[language=C]|n = 25;|}
                \mauvaise{\lstinline[language=C]|n == 25;|}
                \mauvaise{\lstinline[language=C]|25 = n;|}
                \mauvaise{\lstinline[language=C]|n := 25;|}
            \end{reponses}
        \end{multicols}
    \end{question}
}
\element{bank-aff-egal-vs-double-egal}{
    \begin{question}{aff-egal-vs-double-egal-v3}
        Pour affecter la valeur \texttt{7} à la variable \texttt{score}, quelle écriture est correcte ?
        \begin{multicols}{4}
            \begin{reponses}
                \bonne{\lstinline[language=C]|score = 7;|}
                \mauvaise{\lstinline[language=C]|score == 7;|}
                \mauvaise{\lstinline[language=C]|7 = score;|}
                \mauvaise{\lstinline[language=C]|score := 7;|}
            \end{reponses}
        \end{multicols}
    \end{question}
}
\element{affectations}{
    \restituegroupe[1]{bank-aff-egal-vs-double-egal}
}

% ---- aff-point-virgule ----
\element{bank-aff-point-virgule}{
    \begin{question}{aff-point-virgule-v1}
        Quelle instruction est syntaxiquement correcte en C ?
        \begin{multicols}{3}
            \begin{reponses}
                \mauvaise{\lstinline[language=C]|x = x + 1|}
                \bonne{\lstinline[language=C]|x = x + 1;|}
                \mauvaise{\lstinline[language=C]|x = x + 1;;;|}
            \end{reponses}
        \end{multicols}
    \end{question}
}
\element{bank-aff-point-virgule}{
    \begin{question}{aff-point-virgule-v2}
        Quelle instruction est syntaxiquement correcte en C ?
        \begin{multicols}{3}
            \begin{reponses}
                \mauvaise{\lstinline[language=C]|n = n - 2|}
                \bonne{\lstinline[language=C]|n = n - 2;|}
                \mauvaise{\lstinline[language=C]|n = n - 2;;;|}
            \end{reponses}
        \end{multicols}
    \end{question}
}
\element{bank-aff-point-virgule}{
    \begin{question}{aff-point-virgule-v3}
        Quelle instruction est syntaxiquement correcte en C ?
        \begin{multicols}{3}
            \begin{reponses}
                \mauvaise{\lstinline[language=C]|total = total + p|}
                \bonne{\lstinline[language=C]|total = total + p;|}
                \mauvaise{\lstinline[language=C]|total = total + p;;;|}
            \end{reponses}
        \end{multicols}
    \end{question}
}
\element{affectations}{
    \restituegroupe[1]{bank-aff-point-virgule}
}

% ---- aff-incrementation ----
\element{bank-aff-incrementation}{
    \begin{question}{aff-incrementation-v1}
        Après \lstinline[language=C]|int x = 5; x++;|, quelle est la valeur de \texttt{x} ?
        \begin{multicols}{4}
            \begin{reponses}
                \mauvaise{4}
                \bonne{6}
                \mauvaise{5}
                \mauvaise{Erreur de compilation}
            \end{reponses}
        \end{multicols}
    \end{question}
}
\element{bank-aff-incrementation}{
    \begin{question}{aff-incrementation-v2}
        Après \lstinline[language=C]|int x = 9; x++;|, quelle est la valeur de \texttt{x} ?
        \begin{multicols}{4}
            \begin{reponses}
                \mauvaise{8}
                \bonne{10}
                \mauvaise{9}
                \mauvaise{Erreur de compilation}
            \end{reponses}
        \end{multicols}
    \end{question}
}
\element{bank-aff-incrementation}{
    \begin{question}{aff-incrementation-v3}
        Après \lstinline[language=C]|int x = 2; x++;|, quelle est la valeur de \texttt{x} ?
        \begin{multicols}{4}
            \begin{reponses}
                \mauvaise{1}
                \bonne{3}
                \mauvaise{2}
                \mauvaise{Erreur de compilation}
            \end{reponses}
        \end{multicols}
    \end{question}
}
\element{affectations}{
    \restituegroupe[1]{bank-aff-incrementation}
}

% ---- aff-division-entiere ----
% Durci : la division entière positive (arrondi vers le bas) est déjà bien
% connue en fin de séquence. On introduit un opérande négatif : le piège
% classique est de croire que le C arrondit vers le bas (comme une division
% mathématique/floor), alors qu'il tronque vers zéro.
\element{bank-aff-division-entiere}{
    \begin{question}{aff-division-entiere-v1}
        Que vaut \texttt{r} après \lstinline[language=C]|int r = -7 / 2;| ?
        \begin{multicols}{5}
            \begin{reponses}
                \bonne{-3}
                \mauvaise{-4}
                \mauvaise{-3.5}
                \mauvaise{3}
                \mauvaise{-7}
            \end{reponses}
        \end{multicols}
    \end{question}
}
\element{bank-aff-division-entiere}{
    \begin{question}{aff-division-entiere-v2}
        Que vaut \texttt{r} après \lstinline[language=C]|int r = -9 / 2;| ?
        \begin{multicols}{5}
            \begin{reponses}
                \bonne{-4}
                \mauvaise{-5}
                \mauvaise{-4.5}
                \mauvaise{4}
                \mauvaise{-9}
            \end{reponses}
        \end{multicols}
    \end{question}
}
\element{bank-aff-division-entiere}{
    \begin{question}{aff-division-entiere-v3}
        Que vaut \texttt{r} après \lstinline[language=C]|int r = -11 / 4;| ?
        \begin{multicols}{5}
            \begin{reponses}
                \bonne{-2}
                \mauvaise{-3}
                \mauvaise{-2.75}
                \mauvaise{2}
                \mauvaise{-11}
            \end{reponses}
        \end{multicols}
    \end{question}
}
\element{affectations}{
    \restituegroupe[1]{bank-aff-division-entiere}
}

% ---- aff-modulo ----
\element{bank-aff-modulo}{
    \begin{question}{aff-modulo-v1}
        Que vaut \texttt{r} après \texttt{int r = 13 \% 5;} ?
        \begin{multicols}{4}
            \begin{reponses}
                \mauvaise{2.6}
                \mauvaise{2}
                \bonne{3}
                \mauvaise{5}
            \end{reponses}
        \end{multicols}
    \end{question}
}
\element{bank-aff-modulo}{
    \begin{question}{aff-modulo-v2}
        Que vaut \texttt{r} après \texttt{int r = 17 \% 4;} ?
        \begin{multicols}{4}
            \begin{reponses}
                \mauvaise{4.25}
                \mauvaise{4}
                \bonne{1}
                \mauvaise{5}
            \end{reponses}
        \end{multicols}
    \end{question}
}
\element{bank-aff-modulo}{
    \begin{question}{aff-modulo-v3}
        Que vaut \texttt{r} après \texttt{int r = 22 \% 6;} ?
        \begin{multicols}{4}
            \begin{reponses}
                \mauvaise{3.6}
                \mauvaise{3}
                \bonne{4}
                \mauvaise{6}
            \end{reponses}
        \end{multicols}
    \end{question}
}
\element{affectations}{
    \restituegroupe[1]{bank-aff-modulo}
}

% ---- aff-trace (question de trace : tableau structuré au lieu du quadrillage) ----
\element{bank-aff-trace}{
    \begin{question}{aff-trace-v1}
        Pour chacune des lignes suivantes, indiquer la variable modifiée et sa nouvelle valeur.
        \lstinputlisting[language=c,basicstyle=\Large]{sources/codes/affectation-1.c}
        \evaluationProfTableauTrace{2}{%
            \begin{center}
            \begin{tabular}{|c|c|p{4cm}|}
                \hline
                \textbf{Ligne} & \textbf{Variable} & \textbf{Nouvelle valeur} \\
                \hline
                \texttt{val\_i = val\_i + 5;} & & \\[0.3cm]
                \hline
                \texttt{val\_f = val\_i / 2;} & & \\[0.3cm]
                \hline
                \texttt{val\_f = val\_i / 2.0;} & & \\[0.3cm]
                \hline
                \texttt{val\_c = 'C';} & & \\[0.3cm]
                \hline
                \texttt{val\_c++;} & & \\[0.3cm]
                \hline
                \texttt{val\_i = 17 \% 4;} & & \\[0.3cm]
                \hline
                \texttt{val\_i = 20 \% 4;} & & \\[0.3cm]
                \hline
            \end{tabular}
            \end{center}
        }
    \end{question}
}
\element{bank-aff-trace}{
    \begin{question}{aff-trace-v2}
        Pour chacune des lignes suivantes, indiquer la variable modifiée et sa nouvelle valeur.
        \lstinputlisting[language=c,basicstyle=\Large]{sources/codes/affectation-2.c}
        \evaluationProfTableauTrace{2}{%
            \begin{center}
            \begin{tabular}{|c|c|p{4cm}|}
                \hline
                \textbf{Ligne} & \textbf{Variable} & \textbf{Nouvelle valeur} \\
                \hline
                \texttt{val\_i = val\_i + 4;} & & \\[0.3cm]
                \hline
                \texttt{val\_f = val\_i / 2;} & & \\[0.3cm]
                \hline
                \texttt{val\_f = val\_i / 2.0;} & & \\[0.3cm]
                \hline
                \texttt{val\_c = 'D';} & & \\[0.3cm]
                \hline
                \texttt{val\_c++;} & & \\[0.3cm]
                \hline
                \texttt{val\_i = 19 \% 5;} & & \\[0.3cm]
                \hline
                \texttt{val\_i = 22 \% 5;} & & \\[0.3cm]
                \hline
            \end{tabular}
            \end{center}
        }
    \end{question}
}
\element{bank-aff-trace}{
    \begin{question}{aff-trace-v3}
        Pour chacune des lignes suivantes, indiquer la variable modifiée et sa nouvelle valeur.
        \lstinputlisting[language=c,basicstyle=\Large]{sources/codes/affectation-3.c}
        \evaluationProfTableauTrace{2}{%
            \begin{center}
            \begin{tabular}{|c|c|p{4cm}|}
                \hline
                \textbf{Ligne} & \textbf{Variable} & \textbf{Nouvelle valeur} \\
                \hline
                \texttt{val\_i = val\_i + 3;} & & \\[0.3cm]
                \hline
                \texttt{val\_f = val\_i / 2;} & & \\[0.3cm]
                \hline
                \texttt{val\_f = val\_i / 2.0;} & & \\[0.3cm]
                \hline
                \texttt{val\_c = 'E';} & & \\[0.3cm]
                \hline
                \texttt{val\_c++;} & & \\[0.3cm]
                \hline
                \texttt{val\_i = 23 \% 6;} & & \\[0.3cm]
                \hline
                \texttt{val\_i = 25 \% 6;} & & \\[0.3cm]
                \hline
            \end{tabular}
            \end{center}
        }
    \end{question}
}
\element{affectations}{
    \restituegroupe[1]{bank-aff-trace}
}

% ==================== THEME 3 — CONDITIONS ====================

% ---- cond-egalite ----
\element{bank-cond-egalite}{
    \begin{question}{cond-egalite-v1}
        Quel opérateur permet de \textbf{tester} l'égalité entre deux entiers (et non affecter) ?
        \begin{multicols}{4}
            \begin{reponses}
                \mauvaise{\texttt{=}}
                \bonne{\texttt{==}}
                \mauvaise{\texttt{=<}}
                \mauvaise{\texttt{:=}}
            \end{reponses}
        \end{multicols}
    \end{question}
}
\element{bank-cond-egalite}{
    \begin{question}{cond-egalite-v2}
        Quel opérateur permet de \textbf{tester} l'égalité entre deux caractères (et non affecter) ?
        \begin{multicols}{4}
            \begin{reponses}
                \mauvaise{\texttt{=}}
                \bonne{\texttt{==}}
                \mauvaise{\texttt{=<}}
                \mauvaise{\texttt{:=}}
            \end{reponses}
        \end{multicols}
    \end{question}
}
\element{bank-cond-egalite}{
    \begin{question}{cond-egalite-v3}
        Quel opérateur permet de \textbf{tester} l'égalité entre deux variables flottantes (et non affecter) ?
        \begin{multicols}{4}
            \begin{reponses}
                \mauvaise{\texttt{=}}
                \bonne{\texttt{==}}
                \mauvaise{\texttt{=<}}
                \mauvaise{\texttt{:=}}
            \end{reponses}
        \end{multicols}
    \end{question}
}
\element{conditions}{
    \restituegroupe[1]{bank-cond-egalite}
}

% ---- cond-intervalle ----
\element{bank-cond-intervalle}{
    \begin{question}{cond-intervalle-v1}
        Comment écrire la condition \(5 \le x \le 10\) (inclus) ?
        \begin{multicols}{4}
            \begin{reponses}
                \bonne{\lstinline[language=C]|x >= 5 && x <= 10|}
                \mauvaise{\lstinline[language=C]!x > 5 || x < 10!}
                \mauvaise{\lstinline[language=C]|x > 5 && x < 10|}
                \mauvaise{\lstinline[language=C]|x <= 5 && x >= 10|}
            \end{reponses}
        \end{multicols}
    \end{question}
}
\element{bank-cond-intervalle}{
    \begin{question}{cond-intervalle-v2}
        Comment écrire la condition \(0 \le n \le 20\) (inclus) ?
        \begin{multicols}{4}
            \begin{reponses}
                \bonne{\lstinline[language=C]|n >= 0 && n <= 20|}
                \mauvaise{\lstinline[language=C]!n > 0 || n < 20!}
                \mauvaise{\lstinline[language=C]|n > 0 && n < 20|}
                \mauvaise{\lstinline[language=C]|n <= 0 && n >= 20|}
            \end{reponses}
        \end{multicols}
    \end{question}
}
\element{bank-cond-intervalle}{
    \begin{question}{cond-intervalle-v3}
        Comment écrire la condition \(1 \le note \le 6\) (inclus) ?
        \begin{multicols}{4}
            \begin{reponses}
                \bonne{\lstinline[language=C]|note >= 1 && note <= 6|}
                \mauvaise{\lstinline[language=C]!note > 1 || note < 6!}
                \mauvaise{\lstinline[language=C]|note > 1 && note < 6|}
                \mauvaise{\lstinline[language=C]|note <= 1 && note >= 6|}
            \end{reponses}
        \end{multicols}
    \end{question}
}
% v4 : variante supplémentaire demandée explicitement pour l'intervalle
% 5 <= x <= 10 ; la variable est renommée (\texttt{score}) pour ne pas
% dupliquer mot pour mot la v1, qui utilise déjà les mêmes bornes sur \texttt{x}.
\element{bank-cond-intervalle}{
    \begin{question}{cond-intervalle-v4}
        Comment écrire la condition \(5 \le score \le 10\) (inclus) ?
        \begin{multicols}{4}
            \begin{reponses}
                \bonne{\lstinline[language=C]|score >= 5 && score <= 10|}
                \mauvaise{\lstinline[language=C]!score > 5 || score < 10!}
                \mauvaise{\lstinline[language=C]|score > 5 && score < 10|}
                \mauvaise{\lstinline[language=C]|score <= 5 && score >= 10|}
            \end{reponses}
        \end{multicols}
    \end{question}
}
\element{conditions}{
    \restituegroupe[1]{bank-cond-intervalle}
}

% ---- cond-si-erreur-frequente ----
\element{bank-cond-si-erreur-frequente}{
    \begin{question}{cond-si-erreur-frequente-v1}
        Quelle est l'erreur dans le code suivant ?
        \lstinputlisting{sources/codes/cond-erreur-egal.c}
        \begin{multicols}{2}
            \begin{reponses}
                \bonne{\texttt{=} affecte 5 à x au lieu de tester l'égalité (il fallait \texttt{==})}
                \mauvaise{Il manque un point-virgule après le \texttt{if}}
                \mauvaise{Les accolades sont interdites après un \texttt{if}}
                \mauvaise{\texttt{printf} ne peut pas être utilisé dans un \texttt{if}}
            \end{reponses}
        \end{multicols}
    \end{question}
}
\element{bank-cond-si-erreur-frequente}{
    \begin{question}{cond-si-erreur-frequente-v2}
        Quelle est l'erreur dans le code suivant ?
        \lstinputlisting{sources/codes/cond-erreur-egal-2.c}
        \begin{multicols}{2}
            \begin{reponses}
                \bonne{\texttt{=} affecte 10 à note au lieu de tester l'égalité (il fallait \texttt{==})}
                \mauvaise{Il manque un point-virgule après le \texttt{if}}
                \mauvaise{Les accolades sont interdites après un \texttt{if}}
                \mauvaise{\texttt{printf} ne peut pas être utilisé dans un \texttt{if}}
            \end{reponses}
        \end{multicols}
    \end{question}
}
\element{bank-cond-si-erreur-frequente}{
    \begin{question}{cond-si-erreur-frequente-v3}
        Quelle est l'erreur dans le code suivant ?
        \lstinputlisting{sources/codes/cond-erreur-egal-3.c}
        \begin{multicols}{2}
            \begin{reponses}
                \bonne{\texttt{=} affecte 1 à statut au lieu de tester l'égalité (il fallait \texttt{==})}
                \mauvaise{Il manque un point-virgule après le \texttt{if}}
                \mauvaise{Les accolades sont interdites après un \texttt{if}}
                \mauvaise{\texttt{printf} ne peut pas être utilisé dans un \texttt{if}}
            \end{reponses}
        \end{multicols}
    \end{question}
}
\element{conditions}{
    \restituegroupe[1]{bank-cond-si-erreur-frequente}
}

% ---- cond-else-if ----
\element{bank-cond-else-if}{
    \begin{question}{cond-else-if-v1}
        Dans une structure \texttt{if}/\texttt{else if}/\texttt{else}, combien de blocs au maximum peuvent s'exécuter pour un même passage ?
        \begin{multicols}{4}
            \begin{reponses}
                \bonne{1}
                \mauvaise{Autant que de conditions vraies}
                \mauvaise{2}
                \mauvaise{0}
            \end{reponses}
        \end{multicols}
    \end{question}
}
\element{bank-cond-else-if}{
    \begin{question}{cond-else-if-v2}
        Dans une structure \texttt{if}/\texttt{else if}/\texttt{else if}/\texttt{else}, combien de blocs au maximum sont exécutés pour un même passage ?
        \begin{multicols}{4}
            \begin{reponses}
                \bonne{1}
                \mauvaise{Autant que de conditions vraies}
                \mauvaise{2}
                \mauvaise{0}
            \end{reponses}
        \end{multicols}
    \end{question}
}
\element{bank-cond-else-if}{
    \begin{question}{cond-else-if-v3}
        Dans une structure \texttt{if}/\texttt{else}, combien de blocs au maximum peuvent s'exécuter pour un même passage ?
        \begin{multicols}{4}
            \begin{reponses}
                \bonne{1}
                \mauvaise{Autant que de conditions vraies}
                \mauvaise{2}
                \mauvaise{0}
            \end{reponses}
        \end{multicols}
    \end{question}
}
\element{conditions}{
    \restituegroupe[1]{bank-cond-else-if}
}

% ---- cond-et-ou ----
% Durci : un simple "&&" à deux termes se résout par calcul direct sans
% connaître les priorités d'opérateurs. On passe à un mélange && / || à trois
% termes où \&\& est prioritaire sur \textbar\textbar : une lecture naïve
% gauche-à-droite (comme si les deux opérateurs avaient la même priorité)
% donne un résultat différent du résultat correct.
\element{bank-cond-et-ou}{
    \begin{question}{cond-et-ou-v1}
        Que vaut l'expression \lstinline[language=C]!(3 > 1) || (2 > 5) && (4 > 10)! ?
        \begin{multicols}{5}
            \begin{reponses}
                \bonne{Vrai (1)}
                \mauvaise{Faux (0)}
                \mauvaise{Erreur de compilation}
                \mauvaise{0}
                \mauvaise{10}
            \end{reponses}
        \end{multicols}
    \end{question}
}
\element{bank-cond-et-ou}{
    \begin{question}{cond-et-ou-v2}
        Que vaut l'expression \lstinline[language=C]!(5 > 2) || (1 > 9) && (7 > 20)! ?
        \begin{multicols}{5}
            \begin{reponses}
                \bonne{Vrai (1)}
                \mauvaise{Faux (0)}
                \mauvaise{Erreur de compilation}
                \mauvaise{0}
                \mauvaise{20}
            \end{reponses}
        \end{multicols}
    \end{question}
}
\element{bank-cond-et-ou}{
    \begin{question}{cond-et-ou-v3}
        Que vaut l'expression \lstinline[language=C]!(6 > 4) || (0 > 8) && (3 > 15)! ?
        \begin{multicols}{5}
            \begin{reponses}
                \bonne{Vrai (1)}
                \mauvaise{Faux (0)}
                \mauvaise{Erreur de compilation}
                \mauvaise{0}
                \mauvaise{15}
            \end{reponses}
        \end{multicols}
    \end{question}
}
\element{conditions}{
    \restituegroupe[1]{bank-cond-et-ou}
}

% ---- cond-ecrire (écriture de code : reste \evaluationProfGrille, pas un tableau) ----
\element{bank-cond-ecrire}{
    \begin{question}{cond-ecrire-v1}
        Écrire une condition en langage C qui affiche \texttt{"Majeur"} si la variable entière \texttt{age} est supérieure ou égale à 18, et \texttt{"Mineur"} sinon.
        \evaluationProfGrille{2}{9}
    \end{question}
}
\element{bank-cond-ecrire}{
    \begin{question}{cond-ecrire-v2}
        Écrire une condition en langage C qui affiche \texttt{"Positif"} si la variable entière \texttt{n} est strictement supérieure à 0, et \texttt{"Négatif ou nul"} sinon.
        \evaluationProfGrille{2}{9}
    \end{question}
}
\element{bank-cond-ecrire}{
    \begin{question}{cond-ecrire-v3}
        Écrire une condition en langage C qui affiche \texttt{"Admis"} si la variable entière \texttt{note} est supérieure ou égale à 10, et \texttt{"Recalé"} sinon.
        \evaluationProfGrille{2}{9}
    \end{question}
}
\element{conditions}{
    \restituegroupe[1]{bank-cond-ecrire}
}

% ==================== THEME 4 — BOUCLES SIMPLES ====================

% ---- boucle-for-trace ----
\element{bank-boucle-for-trace}{
    \begin{question}{boucle-for-trace-v1}
        Soit le code ci-dessous, qu'affichera le programme en fin d'exécution ?
        \lstinputlisting{sources/codes/boucle-for-1.c}
        \begin{multicols}{6}
            \begin{reponses}
                \bonne{10}
                \mauvaise{15}
                \mauvaise{0}
                \mauvaise{5}
                \mauvaise{4}
                \mauvaise{i}
            \end{reponses}
        \end{multicols}
    \end{question}
}
\element{bank-boucle-for-trace}{
    \begin{question}{boucle-for-trace-v2}
        Soit le code ci-dessous, qu'affichera le programme en fin d'exécution ?
        \lstinputlisting{sources/codes/boucle-for-2.c}
        \begin{multicols}{6}
            \begin{reponses}
                \bonne{6}
                \mauvaise{10}
                \mauvaise{0}
                \mauvaise{4}
                \mauvaise{3}
                \mauvaise{i}
            \end{reponses}
        \end{multicols}
    \end{question}
}
\element{bank-boucle-for-trace}{
    \begin{question}{boucle-for-trace-v3}
        Soit le code ci-dessous, qu'affichera le programme en fin d'exécution ?
        \lstinputlisting{sources/codes/boucle-for-3.c}
        \begin{multicols}{6}
            \begin{reponses}
                \bonne{15}
                \mauvaise{21}
                \mauvaise{0}
                \mauvaise{6}
                \mauvaise{5}
                \mauvaise{i}
            \end{reponses}
        \end{multicols}
    \end{question}
}
\element{boucles}{
    \restituegroupe[1]{bank-boucle-for-trace}
}

% ---- boucle-while-trace ----
\element{bank-boucle-while-trace}{
    \begin{question}{boucle-while-trace-v1}
        Soit le code ci-dessous, qu'affichera le programme en fin d'exécution ?
        \lstinputlisting{sources/codes/boucle-while-1.c}
        \begin{multicols}{5}
            \begin{reponses}
                \bonne{0123}
                \mauvaise{1234}
                \mauvaise{01234}
                \mauvaise{4}
                \mauvaise{Boucle infinie}
            \end{reponses}
        \end{multicols}
    \end{question}
}
\element{bank-boucle-while-trace}{
    \begin{question}{boucle-while-trace-v2}
        Soit le code ci-dessous, qu'affichera le programme en fin d'exécution ?
        \lstinputlisting{sources/codes/boucle-while-2.c}
        \begin{multicols}{5}
            \begin{reponses}
                \bonne{01234}
                \mauvaise{12345}
                \mauvaise{012345}
                \mauvaise{5}
                \mauvaise{Boucle infinie}
            \end{reponses}
        \end{multicols}
    \end{question}
}
\element{bank-boucle-while-trace}{
    \begin{question}{boucle-while-trace-v3}
        Soit le code ci-dessous, qu'affichera le programme en fin d'exécution ?
        \lstinputlisting{sources/codes/boucle-while-3.c}
        \begin{multicols}{5}
            \begin{reponses}
                \bonne{012}
                \mauvaise{123}
                \mauvaise{0123}
                \mauvaise{3}
                \mauvaise{Boucle infinie}
            \end{reponses}
        \end{multicols}
    \end{question}
}
\element{boucles}{
    \restituegroupe[1]{bank-boucle-while-trace}
}

% ---- boucle-for-vs-while ----
\element{bank-boucle-for-vs-while}{
    \begin{question}{boucle-for-vs-while-v1}
        Quelle boucle utiliser lorsque l'on connaît \textbf{à l'avance} le nombre exact de répétitions ?
        \begin{multicols}{3}
            \begin{reponses}
                \bonne{\lstinline[language=C]|for|}
                \mauvaise{\lstinline[language=C]|while| uniquement}
                \mauvaise{\lstinline[language=C]|if|}
            \end{reponses}
        \end{multicols}
    \end{question}
}
\element{bank-boucle-for-vs-while}{
    \begin{question}{boucle-for-vs-while-v2}
        Quelle boucle est la plus adaptée pour répéter un traitement un nombre \textbf{fixe et connu} de fois (par exemple 20 fois) ?
        \begin{multicols}{3}
            \begin{reponses}
                \bonne{\lstinline[language=C]|for|}
                \mauvaise{\lstinline[language=C]|while| uniquement}
                \mauvaise{\lstinline[language=C]|if|}
            \end{reponses}
        \end{multicols}
    \end{question}
}
\element{bank-boucle-for-vs-while}{
    \begin{question}{boucle-for-vs-while-v3}
        Pour parcourir un intervalle dont les bornes sont connues dès le départ, quelle structure de boucle est la plus naturelle en C ?
        \begin{multicols}{3}
            \begin{reponses}
                \bonne{\lstinline[language=C]|for|}
                \mauvaise{\lstinline[language=C]|while| uniquement}
                \mauvaise{\lstinline[language=C]|if|}
            \end{reponses}
        \end{multicols}
    \end{question}
}
\element{boucles}{
    \restituegroupe[1]{bank-boucle-for-vs-while}
}

% ---- boucle-condition-arret ----
\element{bank-boucle-condition-arret}{
    \begin{question}{boucle-condition-arret-v1}
        Quel risque court le programme suivant ?
        \lstinputlisting{sources/codes/boucle-infinie.c}
        \begin{multicols}{2}
            \begin{reponses}
                \bonne{Boucle infinie, car \texttt{i} n'est jamais incrémenté}
                \mauvaise{Erreur de compilation}
                \mauvaise{Le programme n'affiche rien}
                \mauvaise{Le programme s'arrête après un seul tour}
            \end{reponses}
        \end{multicols}
    \end{question}
}
\element{bank-boucle-condition-arret}{
    \begin{question}{boucle-condition-arret-v2}
        Quel risque court le programme suivant ?
        \lstinputlisting{sources/codes/boucle-infinie-2.c}
        \begin{multicols}{2}
            \begin{reponses}
                \bonne{Boucle infinie, car \texttt{n} n'est jamais incrémenté}
                \mauvaise{Erreur de compilation}
                \mauvaise{Le programme n'affiche rien}
                \mauvaise{Le programme s'arrête après un seul tour}
            \end{reponses}
        \end{multicols}
    \end{question}
}
\element{bank-boucle-condition-arret}{
    \begin{question}{boucle-condition-arret-v3}
        Quel risque court le programme suivant ?
        \lstinputlisting{sources/codes/boucle-infinie-3.c}
        \begin{multicols}{2}
            \begin{reponses}
                \bonne{Boucle infinie, car \texttt{k} n'est jamais décrémenté}
                \mauvaise{Erreur de compilation}
                \mauvaise{Le programme n'affiche rien}
                \mauvaise{Le programme s'arrête après un seul tour}
            \end{reponses}
        \end{multicols}
    \end{question}
}
\element{boucles}{
    \restituegroupe[1]{bank-boucle-condition-arret}
}

% ---- boucle-offbyone ----
% Durci : le simple off-by-one sur une borne unique (\texttt{i <= N}) est
% acquis en fin de séquence. On ajoute une seconde condition combinée par
% \texttt{\&\&} qui provoque un arrêt anticipé de la boucle : le piège n'est
% plus de compter les bornes mais de repérer QUELLE condition arrête la
% boucle en premier.
\element{bank-boucle-offbyone}{
    \begin{question}{boucle-offbyone-v1}
        Combien de fois le message \texttt{"Bonjour"} est-il affiché par \lstinline[language=C]|for (int i = 1; i <= 10 && i != 6; i++) { printf("Bonjour"); }| ?
        \begin{multicols}{5}
            \begin{reponses}
                \mauvaise{10}
                \bonne{5}
                \mauvaise{6}
                \mauvaise{4}
                \mauvaise{Boucle infinie}
            \end{reponses}
        \end{multicols}
    \end{question}
}
\element{bank-boucle-offbyone}{
    \begin{question}{boucle-offbyone-v2}
        Combien de fois le message \texttt{"Bonjour"} est-il affiché par \lstinline[language=C]|for (int i = 1; i <= 8 && i != 5; i++) { printf("Bonjour"); }| ?
        \begin{multicols}{5}
            \begin{reponses}
                \mauvaise{8}
                \bonne{4}
                \mauvaise{5}
                \mauvaise{3}
                \mauvaise{Boucle infinie}
            \end{reponses}
        \end{multicols}
    \end{question}
}
\element{bank-boucle-offbyone}{
    \begin{question}{boucle-offbyone-v3}
        Combien de fois le message \texttt{"Bonjour"} est-il affiché par \lstinline[language=C]|for (int i = 1; i <= 12 && i != 4; i++) { printf("Bonjour"); }| ?
        \begin{multicols}{5}
            \begin{reponses}
                \mauvaise{12}
                \bonne{3}
                \mauvaise{4}
                \mauvaise{2}
                \mauvaise{Boucle infinie}
            \end{reponses}
        \end{multicols}
    \end{question}
}
\element{boucles}{
    \restituegroupe[1]{bank-boucle-offbyone}
}

% ---- boucle-ecrire (écriture de code : reste \evaluationProfGrille, pas un tableau) ----
\element{bank-boucle-ecrire}{
    \begin{question}{boucle-ecrire-v1}
        Écrire une boucle en langage C qui affiche les nombres entiers de 1 à 10 (une valeur par ligne).
        \evaluationProfGrille{2}{9}
    \end{question}
}
\element{bank-boucle-ecrire}{
    \begin{question}{boucle-ecrire-v2}
        Écrire une boucle en langage C qui affiche les nombres entiers de 1 à 6 (une valeur par ligne).
        \evaluationProfGrille{2}{9}
    \end{question}
}
\element{bank-boucle-ecrire}{
    \begin{question}{boucle-ecrire-v3}
        Écrire une boucle en langage C qui affiche les nombres entiers de 0 à 8 (une valeur par ligne).
        \evaluationProfGrille{2}{9}
    \end{question}
}
\element{boucles}{
    \restituegroupe[1]{bank-boucle-ecrire}
}

% ==================== THEME 5 — BOUCLES IMBRIQUEES ====================

% ---- imbr-trace-compteur ----
\element{bank-imbr-trace-compteur}{
    \begin{question}{imbr-trace-compteur-v1}
        Soit le code ci-dessous, qu'affichera le programme en fin d'exécution ?
        \lstinputlisting{sources/codes/boucle-imbriquee-1.c}
        \begin{multicols}{6}
            \begin{reponses}
                \mauvaise{3}
                \mauvaise{2}
                \bonne{6}
                \mauvaise{5}
                \mauvaise{9}
                \mauvaise{Boucle infinie}
            \end{reponses}
        \end{multicols}
    \end{question}
}
\element{bank-imbr-trace-compteur}{
    \begin{question}{imbr-trace-compteur-v2}
        Soit le code ci-dessous, qu'affichera le programme en fin d'exécution ?
        \lstinputlisting{sources/codes/boucle-imbriquee-2.c}
        \begin{multicols}{6}
            \begin{reponses}
                \mauvaise{4}
                \mauvaise{2}
                \bonne{8}
                \mauvaise{6}
                \mauvaise{9}
                \mauvaise{Boucle infinie}
            \end{reponses}
        \end{multicols}
    \end{question}
}
\element{bank-imbr-trace-compteur}{
    \begin{question}{imbr-trace-compteur-v3}
        Soit le code ci-dessous, qu'affichera le programme en fin d'exécution ?
        \lstinputlisting{sources/codes/boucle-imbriquee-3.c}
        \begin{multicols}{6}
            \begin{reponses}
                \mauvaise{3}
                \mauvaise{6}
                \bonne{9}
                \mauvaise{12}
                \mauvaise{27}
                \mauvaise{Boucle infinie}
            \end{reponses}
        \end{multicols}
    \end{question}
}
\element{boucles-imbriquees}{
    \restituegroupe[1]{bank-imbr-trace-compteur}
}

% ---- imbr-nb-tours ----
\element{bank-imbr-nb-tours}{
    \begin{question}{imbr-nb-tours-v1}
        Combien de fois le corps de la boucle interne s'exécute-t-il au total dans \lstinline[language=C]|for (int i=0;i<4;i++) { for (int j=0;j<3;j++) { } }| ?
        \begin{multicols}{5}
            \begin{reponses}
                \mauvaise{4}
                \mauvaise{3}
                \mauvaise{7}
                \bonne{12}
                \mauvaise{1}
            \end{reponses}
        \end{multicols}
    \end{question}
}
\element{bank-imbr-nb-tours}{
    \begin{question}{imbr-nb-tours-v2}
        Combien de fois le corps de la boucle interne s'exécute-t-il au total dans \lstinline[language=C]|for (int i=0;i<5;i++) { for (int j=0;j<2;j++) { } }| ?
        \begin{multicols}{5}
            \begin{reponses}
                \mauvaise{5}
                \mauvaise{2}
                \mauvaise{7}
                \bonne{10}
                \mauvaise{1}
            \end{reponses}
        \end{multicols}
    \end{question}
}
\element{bank-imbr-nb-tours}{
    \begin{question}{imbr-nb-tours-v3}
        Combien de fois le corps de la boucle interne s'exécute-t-il au total dans \lstinline[language=C]|for (int i=0;i<2;i++) { for (int j=0;j<6;j++) { } }| ?
        \begin{multicols}{5}
            \begin{reponses}
                \mauvaise{2}
                \mauvaise{6}
                \mauvaise{8}
                \bonne{12}
                \mauvaise{1}
            \end{reponses}
        \end{multicols}
    \end{question}
}
\element{boucles-imbriquees}{
    \restituegroupe[1]{bank-imbr-nb-tours}
}

% ---- imbr-variable-confusion ----
\element{bank-imbr-variable-confusion}{
    \begin{question}{imbr-variable-confusion-v1}
        Dans une boucle imbriquée, pourquoi faut-il utiliser \textbf{deux variables de boucle différentes} (par exemple \texttt{i} et \texttt{j}) ?
        \begin{multicols}{2}
            \begin{reponses}
                \bonne{Réutiliser la même variable pour les deux boucles casse le comptage de la boucle externe}
                \mauvaise{Le langage C interdit d'avoir deux boucles \texttt{for} dans un même programme}
                \mauvaise{Cela n'a aucune importance, les deux variables peuvent avoir le même nom}
                \mauvaise{Une boucle imbriquée n'a besoin que d'une seule variable}
            \end{reponses}
        \end{multicols}
    \end{question}
}
\element{bank-imbr-variable-confusion}{
    \begin{question}{imbr-variable-confusion-v2}
        Dans une boucle imbriquée, pourquoi utilise-t-on des noms de variable de boucle \textbf{distincts} (par exemple \texttt{i} pour la boucle externe et \texttt{k} pour la boucle interne) ?
        \begin{multicols}{2}
            \begin{reponses}
                \bonne{Réutiliser la même variable pour les deux boucles casse le comptage de la boucle externe}
                \mauvaise{Le langage C interdit d'avoir deux boucles \texttt{for} dans un même programme}
                \mauvaise{Cela n'a aucune importance, les deux variables peuvent avoir le même nom}
                \mauvaise{Une boucle imbriquée n'a besoin que d'une seule variable}
            \end{reponses}
        \end{multicols}
    \end{question}
}
\element{bank-imbr-variable-confusion}{
    \begin{question}{imbr-variable-confusion-v3}
        Que se passe-t-il si l'on utilise la \textbf{même variable} \texttt{i} pour piloter à la fois la boucle externe et la boucle interne d'une boucle imbriquée ?
        \begin{multicols}{2}
            \begin{reponses}
                \bonne{La boucle interne modifie \texttt{i}, ce qui casse le comptage de la boucle externe}
                \mauvaise{Le langage C interdit d'avoir deux boucles \texttt{for} dans un même programme}
                \mauvaise{Cela n'a aucune importance, les deux boucles s'exécutent normalement}
                \mauvaise{Une boucle imbriquée n'a besoin que d'une seule variable}
            \end{reponses}
        \end{multicols}
    \end{question}
}
\element{boucles-imbriquees}{
    \restituegroupe[1]{bank-imbr-variable-confusion}
}

% ---- imbr-trace-variable (question de trace : tableau structuré au lieu du quadrillage) ----
\element{bank-imbr-trace-variable}{
    \begin{question}{imbr-trace-variable-v1}
        Pour chaque tour de la boucle externe, indiquer la valeur de \texttt{total} en fin de tour.
        \begin{minipage}[t]{.4\textwidth}
            \lstinputlisting[language=c,basicstyle=\Large]{sources/codes/trace-imbriquee-1.c}
        \end{minipage}
        \begin{minipage}[t]{.59\textwidth}
            \evaluationProfTableauTrace{2}{%
                \begin{center}
                \begin{tabular}{|c|p{4cm}|}
                    \hline
                    \textbf{Tour (valeur de \texttt{i})} & \textbf{Valeur de \texttt{total} en fin de tour} \\
                    \hline
                    1 & \\[0.3cm]
                    \hline
                    2 & \\[0.3cm]
                    \hline
                    3 & \\[0.3cm]
                    \hline
                \end{tabular}
                \end{center}
            }
        \end{minipage}
    \end{question}
}
\element{bank-imbr-trace-variable}{
    \begin{question}{imbr-trace-variable-v2}
        Pour chaque tour de la boucle externe, indiquer la valeur de \texttt{total} en fin de tour.
        \begin{minipage}[t]{.4\textwidth}
            \lstinputlisting[language=c,basicstyle=\Large]{sources/codes/trace-imbriquee-2.c}
        \end{minipage}
        \begin{minipage}[t]{.59\textwidth}
            \evaluationProfTableauTrace{2}{%
                \begin{center}
                \begin{tabular}{|c|p{4cm}|}
                    \hline
                    \textbf{Tour (valeur de \texttt{i})} & \textbf{Valeur de \texttt{total} en fin de tour} \\
                    \hline
                    1 & \\[0.3cm]
                    \hline
                    2 & \\[0.3cm]
                    \hline
                    3 & \\[0.3cm]
                    \hline
                \end{tabular}
                \end{center}
            }
        \end{minipage}
    \end{question}
}
\element{bank-imbr-trace-variable}{
    \begin{question}{imbr-trace-variable-v3}
        Pour chaque tour de la boucle externe, indiquer la valeur de \texttt{total} en fin de tour.
        \begin{minipage}[t]{.4\textwidth}
            \lstinputlisting[language=c,basicstyle=\Large]{sources/codes/trace-imbriquee-3.c}
        \end{minipage}
        \begin{minipage}[t]{.59\textwidth}
            \evaluationProfTableauTrace{2}{%
                \begin{center}
                \begin{tabular}{|c|p{4cm}|}
                    \hline
                    \textbf{Tour (valeur de \texttt{i})} & \textbf{Valeur de \texttt{total} en fin de tour} \\
                    \hline
                    1 & \\[0.3cm]
                    \hline
                    2 & \\[0.3cm]
                    \hline
                    3 & \\[0.3cm]
                    \hline
                    4 & \\[0.3cm]
                    \hline
                \end{tabular}
                \end{center}
            }
        \end{minipage}
    \end{question}
}
\element{boucles-imbriquees}{
    \restituegroupe[1]{bank-imbr-trace-variable}
}



\def\AMCformQuestion#1{{\sc Question #1 :}}
% En mode "catalog" (\ifAMC@catalog, positionné par --mode C), le but est
% UN SEUL exemplaire de chaque question de chaque banque -- pas une copie
% complète par étudiant. \csvreader boucle sur \ficheEtudiants quel que
% soit le mode (--mode C ne change que le comportement interne d'AMCOpen/
% AMCbeginQuestion, pas cette boucle) : sans ce test, le catalogue était
% reproduit une fois par étudiant de la liste (fichier énorme, confirmé
% par l'utilisateur). On appelle donc \sujet une seule fois directement,
% avec des valeurs factices pour \nom/\prenom/\id (normalement définies
% par \csvreader depuis les colonnes du CSV, donc absentes ici).
\makeatletter
\ifAMC@catalog
\makeatother
  \def\nom{}\def\prenom{}\def\id{}%
  \sujet
\else
  % \ficheEtudiants : même indirection que les autres projets AMC gérés par
  % l'interface web (etudiants.tex, \newcommand{\ficheEtudiants}{<nom>.csv}).
  % etudiants.tex fait lui-même le \newcommand : le défaut ci-dessous ne doit
  % être défini QUE quand ce fichier est absent (compilation en place, ici
  % même), sinon \newcommand échoue avec "already defined" et etudiants.tex
  % n'est alors jamais réellement pris en compte.
  \IfFileExists{etudiants.tex}{\input{etudiants.tex}}{\providecommand{\ficheEtudiants}{students.csv}}
  \csvreader[head to column names,]{\ficheEtudiants}{}{\sujet}
\fi
\end{document}
